\chapter{Cloud contre infrastructure sur site}\label{ch:pl:cloud-against-on-premises}

Le 4 novembre 2021, CME Group et Google Cloud ont annoncé un partenariat de dix ans~: CME allait migrer ses technologies vers Google Cloud, à
partir de l'année suivante avec les services de données et de compensation, puis finalement avec l'ensemble de ses marchés. La même année, la plupart des sociétés de trading gardaient chaque chemin d'ordres dans leurs
propres baies, comme le décrit le livre~14. Les deux décisions peuvent être justes, car elles portent sur des charges de travail différentes. Ce chapitre chiffre celle où la
question est la plus souvent posée et la moins souvent calculée --- le cluster de recherche --- avec chaque prix du cloud pris dans une liste publiée et datée,
chaque hypothèse sur le matériel possédé étant énoncée, et le simulateur du chapitre~\ref{ch:pl:distributed-compute-and-schedulers} tournant sur une
capacité interruptible.

\section{Ce qui migre vers le cloud et ce qui ne le peut pas}

Un cloud public (livre~14, chapitre~16) loue des machines à l'heure, dans des régions et des zones de disponibilité que l'entreprise ne contrôle pas. Quatre questions
permettent de classer une charge de travail. Quelle est sa sensibilité à la latence et à la gigue --- le chemin d'ordres, comme le montre le livre~14, se mesure encore dans les baies de l'entreprise, à côté
du moteur d'appariement. Quelle est la variabilité de sa demande --- un balayage qui requiert des milliers de cœurs pendant un week-end et aucun le lundi
relève du cloud. Où se trouvent ses données --- la question de la section suivante. Et qui doit pouvoir les inspecter --- celle de la dernière section.

\begin{definition}[Infrastructure as code]\label{def:pl:cloud-against-on-premises:iac}
L'\emph{infrastructure as code}\index{infrastructure as code} décrit des serveurs, des réseaux, des stockages et leurs permissions dans des
fichiers versionnés à partir desquels un outil les crée, les modifie et les détruit, de sorte qu'un environnement puisse être reconstruit, relu
\end{definition}

L'\omterm{def:pl:cloud-against-on-premises:iac}{infrastructure as code} rend l'élasticité utilisable~: un cluster qui n'existe que pendant qu'un balayage tourne doit être créé et détruit par un
programme, à l'identique à chaque fois, et un cluster possédé décrit de la même façon peut être reconstruit après une panne. Elle s'applique aux deux côtés de
la comparaison de ce chapitre.

\section{L'économie du calcul~: possédé, réservé, à la demande, interruptible}

\begin{definition}[Capacité à la demande, réservée et spot]\label{def:pl:cloud-against-on-premises:capacity}
La \emph{capacité à la demande}\index{capacité à la demande} se loue à la seconde ou à l'heure, sans engagement, au prix public du fournisseur.
La \emph{capacité réservée}\index{capacité réservée} est engagée pour un ou trois ans à un prix horaire inférieur, payé qu'elle soit utilisée ou non.
La \emph{capacité spot}\index{capacité spot} est la capacité inutilisée du fournisseur, louée à un prix inférieur et variable, que le fournisseur
peut reprendre à bref délai, en interrompant ce qui s'exécute dessus.
\end{definition}

Ce chapitre chiffre une instance cloud de 96 threads matériels dans une région des États-Unis, à partir de la liste de prix publique du fournisseur datée du 25 septembre 2026~: 4.284
dollars de l'heure à la demande, 2.834 en réservé pour un an et 1.886 pour trois ans sans paiement initial (upfront), et 1.737 sur le marché spot le 28
septembre. Par thread-heure, cela fait 4.46, 2.95, 1.96 et 1.81 cents. Le matériel possédé n'a pas de liste équivalente. Le modèle construit un serveur à deux sockets de
512 threads à partir du prix de catalogue de son CPU (10\,931~dollars chacun), suppose que les CPU représentent la moitié du prix du serveur, l'amortit sur quatre ans,
l'alimente au prix moyen commercial de l'électricité aux États-Unis (14.53 cents le kilowattheure en juillet 2026) avec une consommation de 1.3~kW et un surcoût de bâtiment
de 1.4, et ajoute 300~\$ par mois d'espace en baie ainsi qu'un deux-centième du salaire annuel d'un ingénieur, soit 250\,000~\$. Ces derniers chiffres sont
des hypothèses, non des faits, et les résultats ci-dessous sont présentés pour des prix de serveur allant jusqu'à trois fois l'hypothèse.

\omcode{code/firm/cloudcost/firm_cloudcost.py}{45}{61}{Le coût mensuel du serveur possédé~: l'achat réparti sur sa durée de vie, l'électricité
avec le surcoût du bâtiment, l'espace et le personnel~; et son coût par thread-heure à pleine utilisation.}\label{lst:pl:cloud-against-on-premises:owned}

Le serveur possédé coûte alors 1\,508~\$ par mois, soit 0.40 cent par thread-heure à pleine utilisation. La comparaison repose sur le taux d'utilisation~: une capacité
payée à chaque heure --- possédée ou réservée --- coûte son prix horaire divisé par la part des heures où elle est utilisée, tandis que la \omterm{def:pl:cloud-against-on-premises:capacity}{capacité à la demande} et la \omterm{def:pl:cloud-against-on-premises:capacity}{capacité spot}
ne coûtent leur prix que lorsqu'elles sont utilisées (\cref{fig:pl:cloud-against-on-premises:curves}).

\begin{omfigure}
\begin{tikzpicture}
\begin{semilogyaxis}[width=11.5cm,height=6cm,xlabel={taux d'utilisation de la capacité payée},ylabel={cents par thread-heure (log)},
  grid=major,grid style={gray!20},tick label style={font=\small},label style={font=\small},table/col sep=comma,
  xmin=0.05,xmax=1,ymin=0.3,ymax=40,xtick={0.1,0.2,0.4,0.6,0.8,1},
  xticklabel={\pgfmathparse{100*\tick}\pgfmathprintnumber[fixed,precision=0]{\pgfmathresult}\%},
  yticklabel={\pgfmathparse{exp(\tick)}\pgfmathprintnumber[fixed,precision=1]{\pgfmathresult}}]
\addplot[omThm,thick] table[x=u,y expr=100*\thisrow{owned}]{figdata/platforms/26-cloud-against-on-premises/curves.csv};
\addplot[omDef,thick,dashed] table[x=u,y expr=100*\thisrow{reserved_3y}]{figdata/platforms/26-cloud-against-on-premises/curves.csv};
\addplot[black,thick] table[x=u,y expr=100*\thisrow{on_demand}]{figdata/platforms/26-cloud-against-on-premises/curves.csv};
\addplot[gray,thick,dashdotted] table[x=u,y expr=100*\thisrow{spot}]{figdata/platforms/26-cloud-against-on-premises/curves.csv};
\draw[omThm,dotted] (axis cs:0.090,0.3) -- (axis cs:0.090,4.46) node[pos=0.25,right,font=\scriptsize]{9~\%};
\draw[omThm,dotted] (axis cs:0.223,0.3) -- (axis cs:0.223,1.81) node[pos=0.1,right,font=\scriptsize]{22~\%};
\node[omThm,font=\scriptsize,anchor=south west] at (axis cs:0.6,0.68) {possédé};
\node[omDef,font=\scriptsize,anchor=south west] at (axis cs:0.75,2.8) {réservé 3 ans};
\node[font=\scriptsize,anchor=south west] at (axis cs:0.6,4.6) {à la demande};
\node[gray,font=\scriptsize,anchor=north west] at (axis cs:0.42,1.75) {spot, sans perte par interruption};
\end{semilogyaxis}
\end{tikzpicture}

\omcaption{Coût par thread-heure effectivement utilisé en fonction du taux d'utilisation de la capacité payée, aux prix cités et pour le serveur
possédé du modèle. La capacité possédée est moins chère que la \omterm{def:pl:cloud-against-on-premises:capacity}{capacité à la demande} au-dessus de 9~\% d'utilisation et moins chère que le spot au-dessus de 22~\%~; la \omterm{def:pl:cloud-against-on-premises:capacity}{capacité réservée} sur trois ans bat
la \omterm{def:pl:cloud-against-on-premises:capacity}{capacité à la demande} au-dessus de 44~\% mais ne bat jamais la possession. Données~: \texttt{fig\_cloudcost.py}.}\label{fig:pl:cloud-against-on-premises:curves}
\end{omfigure}

\begin{table}[htbp]
\centering\small
\begin{tabular}{@{}>{\raggedright\arraybackslash}p{3.8cm}>{\centering\arraybackslash}p{2.5cm}>{\centering\arraybackslash}p{2.5cm}>{\centering\arraybackslash}p{2.5cm}>{\centering\arraybackslash}p{2.5cm}@{}}
\toprule
prix du serveur & par mois & cents par thread-heure & seuil de rentabilité~: à la demande & spot\\
\midrule
tel que supposé (43\,724~\$) & 1\,508~\$ & 0.40 & 9.0~\% & 22.3~\%\\
double & 2\,419~\$ & 0.65 & 14.5~\% & 35.8~\%\\
triple & 3\,330~\$ & 0.89 & 20.0~\% & 49.2~\%\\
\bottomrule
\end{tabular}
\caption{Le taux d'utilisation au-delà duquel le serveur possédé coûte moins cher que la location des mêmes threads, pour trois prix de serveur. Même au triple
du prix supposé, la possession l'emporte pour toute capacité utilisée plus d'un cinquième du temps face à la demande, et la moitié du temps face au spot.}
\label{tab:pl:cloud-against-on-premises:break-even}
\end{table}

Un thread n'est pas une unité de travail~: le CPU possédé et celui de l'instance cloud n'exécutent pas un même backtest à des vitesses différentes, et une société doit comparer
leurs performances aux mesures du chapitre~\ref{ch:pl:accelerators-for-quantitative-work} sur sa propre charge de travail avant de se fier à un seuil de rentabilité.
Un facteur deux en vitesse divise ou multiplie par deux le coût total du serveur possédé par unité de travail, et le seuil avec lui~: 4.5~\% ou
18~\% face à la demande.

\section{Acheter pour une année de demande}

La demande d'un cluster de recherche n'est pas un taux d'utilisation mais un profil. L'année du modèle comprend une base interactive de 1\,500 threads, 1\,000 de plus pendant
les heures ouvrées en semaine, un batch nocturne de 4\,000 de plus en semaine, un balayage de week-end de 12\,000 de plus pendant huit heures chaque samedi, et vingt
pointes imprévues de 8\,000 pendant quatre heures~: une moyenne de 3\,364 threads pour un pic de 13\,500. En classant les heures par demande, on obtient la
courbe de durée de charge, et chaque niveau de capacité est utilisé dans la part des heures où la demande l'atteint.

\omcode{code/firm/cloudcost/firm_cloudcost.py}{88}{105}{Le planificateur de capacité~: la $k$-ième unité de capacité est nécessaire pendant la
part $f(k)$ des heures~; elle est possédée, réservée ou achetée à l'heure, selon ce qui coûte le moins pour cette part.}\label{lst:pl:cloud-against-on-premises:plan}

\begin{omfigure}
\begin{tikzpicture}
\begin{axis}[width=11.5cm,height=5.6cm,xlabel={heures de l'année, classées par demande},ylabel={threads},
  grid=major,grid style={gray!20},tick label style={font=\small},label style={font=\small},table/col sep=comma,
  xmin=0,xmax=8760,ymin=0,ymax=15000,scaled y ticks=false,yticklabel style={/pgf/number format/1000 sep={\,}},
  xticklabel style={/pgf/number format/1000 sep={\,}}]
\addplot[omDef,thick,const plot] table[x=hours,y=threads]{figdata/platforms/26-cloud-against-on-premises/duration.csv};
\fill[omThm!15] (axis cs:0,0) rectangle (axis cs:8760,5500);
\addplot[omDef,thick,const plot] table[x=hours,y=threads]{figdata/platforms/26-cloud-against-on-premises/duration.csv};
\draw[omThm,dashed] (axis cs:788,0) -- (axis cs:788,15000) node[pos=0.93,right,font=\scriptsize]{9~\% des heures};
\node[omThm,font=\scriptsize] at (axis cs:5000,3000) {possédé~: 5\,500 threads (environ 11 serveurs)};
\node[font=\scriptsize,anchor=south west] at (axis cs:900,9500) {acheté à l'heure~: 496 heures au-dessus de 5\,500};
\end{axis}
\end{tikzpicture}

\omcaption{La courbe de durée de charge de l'année modèle. Tout niveau de demande atteint pendant plus de 9~\% des heures --- le seuil de rentabilité face
à la demande --- coûte moins cher en propriété~: ici jusqu'aux 5\,500 threads du batch nocturne des jours ouvrés. Les balayages du week-end et les rafales imprévues au-dessus
occupent 496 heures et sont achetés à l'heure. Données~: \texttt{fig\_cloudcost.py}.}\label{fig:pl:cloud-against-on-premises:duration}
\end{omfigure}

\begin{omfigure}
\begin{tikzpicture}
\begin{axis}[width=11.5cm,height=5.4cm,xbar stacked,bar width=11pt,xmin=0,xmax=1500,y dir=reverse,
  ytick=data,yticklabels={à la demande seule,réservé et à la demande,possédé et à la demande,possédé et rafales spot},
  xlabel={coût annuel (milliers de dollars)},grid=major,grid style={gray!20},tick label style={font=\small},label style={font=\small},
  table/col sep=comma,legend style={font=\small,at={(0.5,1.03)},anchor=south,legend columns=3,draw=none,
  /tikz/every even column/.append style={column sep=3mm}}]
\addplot[fill=omThm!60,draw=omThm] table[y expr=\coordindex,x expr=\thisrow{owned}/1000]{figdata/platforms/26-cloud-against-on-premises/plans.csv};
\addplot[fill=omDef!40,draw=omDef] table[y expr=\coordindex,x expr=\thisrow{reserved}/1000]{figdata/platforms/26-cloud-against-on-premises/plans.csv};
\addplot[fill=gray!40,draw=gray] table[y expr=\coordindex,x expr=\thisrow{metered}/1000]{figdata/platforms/26-cloud-against-on-premises/plans.csv};
\legend{possédé,réservé (3 ans),à l'heure}
\node[font=\scriptsize,anchor=west] at (axis cs:1320,0) {1\,315};
\node[font=\scriptsize,anchor=west] at (axis cs:945,1) {940};
\node[font=\scriptsize,anchor=west] at (axis cs:366,2) {361};
\node[font=\scriptsize,anchor=west] at (axis cs:287,3) {282};
\end{axis}
\end{tikzpicture}

\omcaption{La demande de l'année modèle, achetée de quatre façons. À la demande seule, le coût atteint 1.32~million de dollars~; réserver les 2\,500 threads utilisés la plupart du temps
ramène le coût à 0.94~million de dollars~; posséder 5\,500 threads et acheter le reste à la demande, à 0.36~million de dollars~; faire tourner 80~\% des rafales sur spot, à
0.28~million de dollars. Données~: \texttt{fig\_cloudcost.py}.}\label{fig:pl:cloud-against-on-premises:plans}
\end{omfigure}

\begin{definition}[Débordement vers le cloud]\label{def:pl:cloud-against-on-premises:bursting}
Le \emph{débordement vers le cloud}\index{débordement vers le cloud} (cloud bursting) exécute une charge de travail sur les capacités propres de l'entreprise jusqu'à leur taille maximale et envoie l'excédent vers une
capacité cloud louée, de sorte que le cluster possédé soit dimensionné pour la charge fréquente et que le cloud absorbe les pics rares.
\end{definition}

La solution optimale de la \cref{fig:pl:cloud-against-on-premises:plans} est hybride, et elle l'est à cause de la forme de la demande, et non à cause d'un
prix particulier~: la base et le batch nocturne sont utilisés plus de 9~\% du temps et coûtent moins cher en propriété~; les balayages et rafales du week-end, 496 heures
de l'année, coûtent moins cher en location. Une entreprise dont la recherche tourne à plein régime chaque nuit possède ses capacités~; une entreprise dont la recherche consiste en quelques énormes balayages par an
les loue~; la plupart se situent entre les deux.

\section{Capacité spot et coût de l'interruption}

La \omterm{def:pl:cloud-against-on-premises:capacity}{capacité spot} est moins chère de 59~\% qu'à la demande pour cette instance, et le conseiller du fournisseur la place dans la bande la plus élevée de fréquence
d'interruption, au-dessus de 20~\% des instances récupérées au cours du dernier mois --- soit au moins 0.03~\% par heure. Le balayage de 10\,000 backtests
du chapitre~\ref{ch:pl:distributed-compute-and-schedulers}, dont quarante durent sept heures et demie chacun, tourne sur 24 instances, chacune facturée jusqu'à la fin de sa dernière tâche.

\omcode{code/firm/cloudcost/firm_cloudcost.py}{109}{124}{Checkpointing sous forme de graphe de tâches~: chaque backtest découpé en morceaux enchaînés, de sorte qu'une
interruption ne fasse perdre au plus qu'un morceau.}\label{lst:pl:cloud-against-on-premises:ckpt}

\begin{omfigure}
\begin{tikzpicture}
\begin{semilogxaxis}[width=11.5cm,height=5.6cm,xlabel={interruptions par heure d'instance (échelle logarithmique)},ylabel={coût du balayage (\$)},
  grid=major,grid style={gray!20},tick label style={font=\small},label style={font=\small},table/col sep=comma,
  xmin=0.0002,xmax=0.3,ymin=0,ymax=220]
\addplot[omThm,thick,mark=*] table[x=rate_per_hour,y=cost_none]{figdata/platforms/26-cloud-against-on-premises/burst.csv};
\addplot[omDef,thick,mark=square*] table[x=rate_per_hour,y=cost_ckpt]{figdata/platforms/26-cloud-against-on-premises/burst.csv};
\draw[black,dashed] (axis cs:0.0002,202.79) -- (axis cs:0.3,202.79) node[pos=0.1,below right,font=\scriptsize]{à la demande~: 203~\$, 7.5 heures};
\draw[gray,dotted] (axis cs:0.000306,0) -- (axis cs:0.000306,220) node[pos=0.55,right,font=\scriptsize,align=left]{bande la plus élevée\\du conseiller};
\node[omThm,font=\scriptsize,anchor=south east] at (axis cs:0.2,150) {spot, reprise à zéro};
\node[omThm,font=\scriptsize,anchor=east] at (axis cs:0.17,184) {59 heures};
\node[omDef,font=\scriptsize,anchor=north west,align=left] at (axis cs:0.0005,72) {spot, checkpoints toutes les 15 minutes~:\\80--82~\$, moins de 9 heures};
\end{semilogxaxis}
\end{tikzpicture}

\omcaption{Le balayage du chapitre~13 sur 24 instances spot, face au taux d'interruption (échelle logarithmique~; chaque instance facturée jusqu'à la fin de sa dernière tâche).
Sans checkpoints, les interruptions relancent les longs backtests depuis le début~: à la bande du conseiller, le balayage coûte 94~\$ et dure 14.4
heures au lieu de 7.5~; à 0.2 par heure, 185~\$ et 59 heures. Avec des checkpoints toutes les 15 minutes, il coûte 80 à 82~\$ et se termine en moins de 9 heures
à chaque taux. Données~: \texttt{fig\_cloudcost.py}.}\label{fig:pl:cloud-against-on-premises:burst}
\end{omfigure}

À la demande, le balayage coûte 203~\$ et dure 7.5 heures, la durée de son backtest le plus long. Sur spot sans interruption, il coûte 82~\$. Ce
qu'une interruption coûte, c'est surtout du temps, et cela touche les longues tâches (la \cref{fig:pl:cloud-against-on-premises:burst})~: un backtest de 7.5 heures
interrompu après six heures repart de zéro, si bien qu'une seule interruption de l'une des quarante longues tâches allonge presque du double la durée du balayage. Découper
chaque backtest en morceaux de 15 minutes qui écrivent un point de reprise --- le point de reprise d'entraînement (checkpoint) du livre~12, chapitre~23, appliqué aux backtests --- borne la
perte à un seul morceau, et le balayage coûte le même prix à chaque taux d'interruption que montre le graphique. La \omterm{def:pl:cloud-against-on-premises:capacity}{capacité spot} est réservée aux travaux qui peuvent être interrompus,
et le checkpointing est ce qui rend un travail interruptible.

\section{Gravité des données et frais de transfert de données}

\begin{definition}[Gravité des données]\label{def:pl:cloud-against-on-premises:gravity}
La \emph{gravité des données}\index{gravité des données} est la tendance du calcul à se déplacer vers l'endroit où se trouvent déjà de grands volumes de données, parce que déplacer les données coûte
plus de temps et de frais de transfert de données que déplacer le calcul.
\end{definition}

Les frais de transfert de données du fournisseur (livre~14, chapitre~16) sont asymétriques~: les données qui entrent dans la région sont gratuites, celles qui en sortent vers internet
coûtent 9 centimes de dollar le gigaoctet pour les 10 premiers téraoctets du mois, puis un tarif dégressif par paliers jusqu'à 5 centimes de dollar au-delà de 150 téraoctets. Conserver une \omterm{def:pl:a-tick-store-on-open-formats:store}{base de ticks} de 500 téraoctets
dans le stockage objet standard du cloud (au prix du chapitre~\ref{ch:pl:capturing-and-storing-tick-data}) coûte 11\,776~\$ par mois~; en extraire
une copie complète une seule fois coûte 29\,491~\$, soit environ vingt mois d'un serveur possédé. Le lieu de la \omterm{def:pl:a-tick-store-on-open-formats:store}{base de ticks} décide donc de l'endroit où tourne la recherche~: une fois
qu'elle est dans le cloud, la recherche la suit, et ce sont les résultats --- quelques gigaoctets par balayage --- qui doivent franchir la frontière. Une architecture hybride conserve
une copie faisant autorité et envoie le calcul vers elle, ou conserve deux copies et paie une fois pour remplir la seconde.

\section{Sécurité, contrôle et régulateurs}

\begin{definition}[Modèle de responsabilité partagée]\label{def:pl:cloud-against-on-premises:shared}
Le \emph{modèle de responsabilité partagée}\index{modèle de responsabilité partagée} répartit la sécurité entre un fournisseur cloud, responsable de
l'infrastructure qui fait tourner ses services, et le client, responsable de ce qu'il construit et configure dessus --- données, identités, règles
réseau, chiffrement --- dans une mesure qui dépend des services qu'il choisit.
\end{definition}

La plupart des incidents cloud dans les établissements financiers se situent du côté client de cette ligne~: un bucket de stockage laissé public, une clé versionnée dans un dépôt,
une permission accordée trop largement --- les sujets du chapitre~\ref{ch:pl:security}. Les régulateurs ajoutent une seconde ligne. Dans l'Union européenne, les
orientations de l'EBA sur l'externalisation, applicables depuis le 30 septembre 2019 et qui ont intégré sa recommandation antérieure sur l'externalisation vers le cloud, exigent une
stratégie de sortie documentée pour toute fonction critique ou importante externalisée~: un plan, testé, permettant de transférer la fonction ailleurs si le fournisseur
défaille ou si la relation prend fin. Pour un cluster de recherche, c'est facile~: la moitié possédée d'une architecture hybride constitue la sortie. Pour une voie d'ordres ou un système de
compensation, c'est une seconde implémentation complète, ce qui explique en partie pourquoi un passage comme celui de CME est un chantier de dix ans. La concentration --- de nombreuses entreprises qui dépendent
des mêmes quelques fournisseurs --- est le risque opérationnel (livre~6, chapitre~28) que les superviseurs surveillent le plus.

\begin{dated}{2026-09}{Prix du cloud et règles applicables}\label{dat:pl:cloud-against-on-premises:prices}
La liste de prix publique d'AWS du 25 septembre 2026 donne, pour une instance c7i.24xlarge (96 vCPU) dans la région US East (N. Virginia) sous Linux, 4.284~\$
l'heure à la demande, 2.83387 en réservation pour un an et 1.88608 pour trois ans, sans paiement initial (upfront)~; son flux de prix spot indiquait 1.7372 le 28
septembre, et son Spot Instance Advisor plaçait la fréquence d'interruption de l'instance dans la bande la plus élevée (au-dessus de 20~\% dans le dernier mois), avec
65~\% d'économies par rapport à la demande. Le transfert de données sortant de la région vers internet coûte 0.09~\$ par gigaoctet pour les 10~To du mois, 0.085
pour les 40 suivants, 0.07 pour les 100 suivants et 0.05 au-delà de 150~To~; l'entrée est gratuite. AMD affiche son EPYC 9755 à 128 cœurs à 10\,931~\$ (prix
pour 1\,000 unités)~; l'EIA indique un prix moyen de l'électricité commerciale aux États-Unis de 14.53 centimes de dollar par kWh pour juillet 2026. CME Group et Google Cloud ont annoncé
leur partenariat de dix ans, avec un investissement en capital d'un milliard de dollars de la part de Google, le 4 novembre 2021. Les orientations de l'EBA sur l'externalisation
(EBA/GL/2019/02) s'appliquent depuis le 30 septembre 2019.
\end{dated}

\section{Tutoriel~: posséder, louer ou les deux}\label{tut:pl:cloud-against-on-premises:tut}

\begin{tutorial}
\textbf{Objectif.} Chiffrer l'année d'un cluster de recherche de toutes les façons, trouver le taux d'utilisation du seuil de rentabilité et lancer un balayage sur la \omterm{def:pl:cloud-against-on-premises:capacity}{capacité spot}.
\textbf{État final~:} \cref{fig:pl:cloud-against-on-premises:curves,fig:pl:cloud-against-on-premises:plans,fig:pl:cloud-against-on-premises:burst}.
\begin{tutsteps}
\item \textbf{Prix}~: \texttt{pl\_cloudcost.PRICES} et \texttt{SERVER}~; \texttt{per\_thread()} et \texttt{break\_evens(k)} pour
$k=1,2,3$.
\item \textbf{Demande}~: \texttt{demand()}~; tracer sa courbe de durée de charge.
\item \textbf{Plans}~: \texttt{plans()}, les quatre façons d'acheter l'année.
\item \textbf{Débordement}~: \texttt{burst(rate, price, checkpoint)} selon les taux d'interruption, avec et sans checkpoints.
\item \textbf{Données}~: \texttt{firm\_cloudcost.transfer\_out} pour la \omterm{def:pl:a-tick-store-on-open-formats:store}{base de ticks} et pour les résultats d'un balayage.
\end{tutsteps}
\textbf{Pour aller plus loin.} Remplacer le thread comme unité par des backtests mesurés par heure sur chaque machine~; laisser la \omterm{def:pl:cloud-against-on-premises:capacity}{capacité réservée} couvrir le
batch nocturne lorsque la société ne peut pas acheter de serveurs à temps.
\end{tutorial}

\section{Construction~: le modèle de coût du cloud}\label{bld:pl:cloud-against-on-premises:cloudcost}

\begin{build}
\bfield{Rôle} Chiffrer tout plan de calcul à partir de listes datées et citées et d'hypothèses énoncées, et choisir le mélange de capacité possédée, réservée, à la demande et
\omterm{def:pl:cloud-against-on-premises:capacity}{capacité spot} pour un profil de demande.
\bfield{Interface} \texttt{PriceItem}, \texttt{OwnedServer}, \texttt{cost\_per\_used\_hour}, \texttt{break\_even}, \texttt{plan\_capacity},
\texttt{checkpointed}, \texttt{billed\_node\_hours}, \texttt{transfer\_out}.
\bfield{Règles} Chaque prix porte sa source et sa date~; chaque hypothèse est nommée et accompagnée d'une analyse de sensibilité~; les plans sont comparés sur le
même profil de demande~; la capacité interruptible n'est réservée qu'aux travaux dotés de checkpoints~; les frais de transfert de données sont comptés dans les deux sens.
\bfield{Tests d'acceptation} \texttt{code/firm/cloudcost/tests/}~: le coût mensuel et horaire du serveur possédé~; les coûts fixes et à l'usage et le
seuil de rentabilité~; le planificateur sur une demande à deux niveaux, avec \omterm{def:pl:cloud-against-on-premises:capacity}{capacité réservée} et spot~; les tronçons de checkpoints, leur enchaînement et la facturation~; les paliers des
frais de transfert de données.
\bfield{Pour aller plus loin} Savings plans et remises d'engagement d'utilisation~; taux d'interruption issus d'un historique mesuré~; niveaux de stockage et frais de récupération dans
la comparaison sur la \omterm{def:pl:cloud-against-on-premises:gravity}{gravité des données}.
\end{build}

\begin{omsources}
\item CME Group et Google Cloud, communiqué de presse du 4 novembre 2021.
\item Listes de prix publiques d'AWS (EC2 à la demande et réservées, transfert de données), flux de prix spot et Spot Instance Advisor~; AWS, \emph{Shared
Responsibility Model}.
\item AMD, page produit EPYC 9755~; U.S. Energy Information Administration, \emph{Electric Power Monthly}, tableau~5.6.A.
\item European Banking Authority, \emph{Guidelines on outsourcing arrangements}, EBA/GL/2019/02.
\item Le livre~14, chapitre~16 (cloud public, zones de disponibilité, frais de transfert de données)~; le livre~16, chapitre~20 (coût total de possession).
\end{omsources}

\section{Exercices}

\begin{exercise}[$\star$]\label{exo:pl:cloud-against-on-premises:1}
Combien coûte une heure de thread à la demande, en réservation sur trois ans et en spot, aux prix cités~?
\end{exercise}

\begin{exercise}[$\star$]\label{exo:pl:cloud-against-on-premises:2}
À partir de quel taux d'utilisation la \omterm{def:pl:cloud-against-on-premises:capacity}{capacité réservée} sur trois ans est-elle moins chère que la \omterm{def:pl:cloud-against-on-premises:capacity}{capacité à la demande}~?
\end{exercise}

\begin{exercise}[$\star$]\label{exo:pl:cloud-against-on-premises:3}
Combien coûte-t-il de sortir de la région, une fois, une \omterm{def:pl:a-tick-store-on-open-formats:store}{base de ticks} de 500~To~?
\end{exercise}

\begin{exercise}[$\star\star$]\label{exo:pl:cloud-against-on-premises:4}
Pourquoi le planificateur ne choisit-il jamais la \omterm{def:pl:cloud-against-on-premises:capacity}{capacité réservée} lorsque la possession des serveurs est autorisée~?
\end{exercise}

\begin{exercise}[$\star\star$]\label{exo:pl:cloud-against-on-premises:5}
Pourquoi une seule interruption d'un long backtest double-t-elle presque la durée du balayage~?
\end{exercise}

\begin{exercise}[$\star\star$]\label{exo:pl:cloud-against-on-premises:6}
Les processeurs de la société exécutent un backtest deux fois plus vite que ceux de l'instance cloud. Qu'advient-il du seuil de rentabilité~?
\end{exercise}

\begin{exercise}[$\star\star\star$]\label{exo:pl:cloud-against-on-premises:7}
\emph{Programmation.} Relancer \texttt{plans()} sur \texttt{demand(spread\_sweep=True)}, où les heures de thread du balayage du week-end s'exécutent dans les cinq
\omterm{def:pl:the-risk-grid:grid}{fenêtres de batch} en semaine à la place. Quelle capacité le planificateur possède-t-il, et combien coûte l'année~?
\end{exercise}

\begin{exercise}[$\star\star\star$]\label{exo:pl:cloud-against-on-premises:8}
\emph{Trouver la faille.} «~Le spot coûte 59~\% de moins, donc on fera tourner tout le cluster de recherche sur spot.~»
\end{exercise}

\section{Problème~: posséder, louer ou les deux}

\begin{problem}[{Problème du week-end --- posséder, louer ou les deux}]\label{pb:pl:cloud-against-on-premises:1}
Les prix du chapitre, le serveur possédé, le profil de demande et le balayage.

\medskip\noindent\textbf{Partie I --- Prix.}
\begin{enumerate}
\item Quels sont les quatre prix du cloud par heure de thread, et d'où viennent-ils~?
\item Qu'est-ce qui est cité et qu'est-ce qui est supposé dans le coût du serveur possédé~?
\item Combien coûte le serveur possédé par mois et par heure de thread~?
\item Comment le taux d'utilisation entre-t-il dans la comparaison~?
\item Quels sont les seuils de rentabilité par rapport à la demande et au spot, et comment évoluent-ils avec le prix du serveur~?
\end{enumerate}

\medskip\noindent\textbf{Partie II --- L'année.}
\begin{enumerate}[resume]
\item Quel est le profil de demande du modèle, sa moyenne et son pic~?
\item Comment le planificateur décide-t-il chaque unité de capacité~?
\item Quelle capacité possède-t-il, et combien d'heures sont achetées à l'heure~?
\item Combien coûte chacun des quatre plans~?
\item Pourquoi l'optimum est-il un hybride~?
\end{enumerate}

\medskip\noindent\textbf{Partie III --- Spot et données.}
\begin{enumerate}[resume]
\item Que dit le fournisseur des interruptions de cette instance~?
\item Quels sont le coût et la durée du balayage en spot sans checkpoints~?
\item Que changent les checkpoints~?
\item Combien coûte le stockage de la \omterm{def:pl:a-tick-store-on-open-formats:store}{base de ticks} dans le cloud, et combien coûte sa sortie~?
\item Qu'exige le régulateur avant qu'une fonction critique ne migre vers un cloud~?
\end{enumerate}

\medskip\noindent\textbf{Partie IV --- Le verdict.}
\begin{enumerate}[resume]
\item Énoncer le \emph{résultat nommé}~: le taux d'utilisation au-delà duquel posséder un cluster de recherche coûte moins cher que le louer aux prix cités, et
le coût du balayage du chapitre~13 sur \omterm{def:pl:cloud-against-on-premises:capacity}{capacité spot} avec préemption, comparé à la demande.
\item Laquelle des hypothèses du modèle vérifierait-on en premier, et comment~?
\item Où la \omterm{def:pl:a-tick-store-on-open-formats:store}{base de ticks} doit-elle être hébergée~?
\item Qu'est-ce qui conduirait à migrer tout le cluster vers le cloud~?
\item En une phrase~: qu'est-ce qui tranche entre posséder et louer le calcul~?
\end{enumerate}
\end{problem}

\section{Questions d'entretien}

\begin{interviewq}[$\star$ \iqroles{développeur}]\label{iq:pl:cloud-against-on-premises:1}
Quelles sont les instances à la demande, réservées et spot, et quand utiliseriez-vous chacune~?
\end{interviewq}

\begin{interviewq}[$\star\star$ \iqroles{développeur}]\label{iq:pl:cloud-against-on-premises:2}
Comment décideriez-vous s'il faut acheter un cluster de recherche ou en louer un~?
\end{interviewq}

\begin{interviewq}[$\star\star$ \iqroles{développeur}]\label{iq:pl:cloud-against-on-premises:3}
Comment lancez-vous un grand balayage de backtests sur une capacité qui peut être retirée à tout moment~?
\end{interviewq}

\begin{interviewq}[$\star\star$ \iqroles{développeur}]\label{iq:pl:cloud-against-on-premises:4}
Votre \omterm{def:pl:a-tick-store-on-open-formats:store}{base de ticks} fait 500~To sur site. Une équipe veut faire tourner la recherche dans le cloud. Que lui dites-vous~?
\end{interviewq}

\begin{interviewq}[$\star\star\star$ \iqroles{développeur}]\label{iq:pl:cloud-against-on-premises:5}
Concevez une plateforme de recherche hybride qui utilise une capacité possédée pour la charge de base et le cloud pour les pics.
\end{interviewq}
